% =====================================================================
%  9. Brecha de investigación
% =====================================================================
\section{Brecha de investigación y aporte que se pretende estudiar}
\label{sec:brecha}

\subsection{Qué está establecido y qué no}
\label{sec:brecha-balance}

La tabla~\ref{tab:balance} resume, por línea, lo que la literatura permite dar por
establecido y lo que queda sin responder para el problema de esta investigación.

\begin{xltabular}{\textwidth}{@{}c L{5.1cm} X@{}}
\caption{Balance por línea: establecido frente a no resuelto}
\label{tab:balance} \\
\toprule
\textbf{L.} & \textbf{Lo que está establecido} & \textbf{Lo que no está resuelto para este problema} \\
\midrule
\endfirsthead
\multicolumn{3}{@{}l}{\small\textit{(continuación de la tabla \ref{tab:balance})}} \\
\toprule
\textbf{L.} & \textbf{Lo que está establecido} & \textbf{Lo que no está resuelto para este problema} \\
\midrule
\endhead
\bottomrule
\endfoot

A &
La calidad de un recurso educativo es multidimensional y la exactitud del contenido
es su dimensión dominante; existen instrumentos validados de revisión humana. &
Los instrumentos evalúan el recurso aislado. No existe instrumento para la
consistencia \emph{entre} documentos de una misma colección. \\

A &
El aseguramiento de calidad en repositorios combina política, tecnología y
participación comunitaria. &
Lo automatizado se detiene en metadatos, enlaces y recomendaciones; ningún
instrumento del inventario verifica el contenido. \\

B &
La IA generativa redistribuye la carga docente; la revisión de sus salidas es trabajo
nuevo. &
No hay evidencia sobre el efecto de un asistente de curación documental sobre el
tiempo y la carga percibida del docente. \\

B &
La curación de materiales es un tema marginal en las síntesis de AIED
(4,5\,\% y 7,6\,\% de las revisiones). &
No se sabe si el escaso volumen refleja ausencia de investigación primaria o ausencia
de síntesis. \\

C &
La tarea de consistencia se formaliza con tres clases y evidencia localizada; la
contradicción es la clase más difícil y la neutral la peor reconocida. &
Los resultados provienen de contratos y texto enciclopédico en inglés. No hay
evidencia sobre documentación de cursos ni sobre corpus en español. \\

C &
Las métricas basadas en inferencia superan a las de n-gramas para consistencia
factual; la granularidad atómica alinea mejor con el juicio humano. &
Las métricas se degradan con fuentes de más de 200 \emph{tokens} y fallan ante
inconsistencias sutiles. No hay protocolo establecido para documentos largos de
curso. \\

D &
Ningún recuperador ni \emph{embedding} es universalmente superior; la combinación de
recuperación léxica y densa ayuda en algunos casos. &
La selección para un corpus en español de dominio académico no puede justificarse con
los bancos de prueba disponibles, que son en inglés. \\

D &
El patrón RAG ancla la generación en documentos y permite actualizar el conocimiento
cambiando el índice. &
La formulación fundacional \textbf{asume un índice razonablemente consistente} y no
aborda qué ocurre cuando la fuente se contradice internamente. \\

D$'$ &
Existen tres marcos de evaluación de RAG con concordancia humana reportada y
diagnóstico por componente. &
Sus jueces pierden concordancia al cambiar de idioma ($\tau$ de Kendall 0,33 al pasar
al español). No pueden adoptarse sin validación propia. \\

E &
El patrón de razonamiento y acción produce trayectorias más ancladas en evidencia;
sus modos de fallo están caracterizados. &
No está demostrado que un agente supere a una canalización fija en tareas
estructuradas como la curación documental. \\

E &
Los agentes con herramientas son vulnerables a instrucciones inyectadas en el
contenido que procesan, sobre todo en campos de texto libre. &
No hay evaluación de esa superficie de ataque en sistemas que ingieren documentos
educativos cargados por usuarios. \\

F &
Las redes modernas están descalibradas y tienden a la sobreconfianza; la calibración
exige un conjunto independiente con etiquetas verdaderas. &
No se sabe si un puntaje de confianza construido heurísticamente sobre señales de
recuperación y detección guarda relación con la corrección observada. \\

F &
Los puntajes de confianza mal calibrados pueden inducir confianza inapropiada; solo
el 14\,\% de los estudios considera el continuo completo. &
No hay evidencia sobre cómo un docente calibra su confianza frente a sugerencias de
curación acompañadas de evidencia. \\

G &
La memoria de retroalimentación permite aprovechar decisiones previas sin
reentrenar; sus riesgos de recolección están documentados. &
No se ha medido si los precedentes de decisiones docentes mejoran la concordancia con
decisiones futuras sobre instancias no usadas para construir el contexto. \\

G &
Existen modelo de procedencia y marcos de gobernanza que fijan requisitos de
trazabilidad y supervisión. &
Los marcos prescriben qué documentar; no establecen si el registro resultante es
suficiente para que un docente reconstruya y audite una decisión. \\

H &
El repertorio metodológico de evaluación de artefactos está consolidado; existe un
instrumento validado de carga percibida. &
El campo lo aplica de forma irregular: 65\,\% de las revisiones de calidad crítica a
media y 51,5\,\% sin acuerdo entre revisores reportado. \\
\end{xltabular}

\subsection{Formulación de la brecha}
\label{sec:brecha-formulacion}

Las filas anteriores convergen en un punto. Las dos literaturas que tocan el problema
lo hacen desde supuestos incompatibles con él:

\begin{itemize}
  \item La literatura de \textbf{calidad de recursos educativos} sabe evaluar el
        recurso, pero su automatización no llega al contenido y su unidad de análisis
        es el documento aislado, no la relación entre documentos.
  \item La literatura de \textbf{recuperación aumentada e inferencia documental} sabe
        verificar afirmaciones contra evidencia, pero \textbf{asume que el corpus de
        referencia es consistente}. \textcite{lewis2020retrieval} no abordan qué
        ocurre cuando la fuente se contradice internamente, y
        \textcite{gao2024retrieval} advierten que la información errónea recuperada
        puede ser peor que la ausencia de información.
\end{itemize}

\begin{estado}[colframe=ecAzul]
\small
\textbf{Brecha.} No existe evidencia sobre el comportamiento de un sistema de
recuperación aumentada con supervisión humana cuando \textbf{la inconsistencia del
propio corpus es el objeto de detección y no un supuesto violado}, y menos aún sobre
documentación de cursos en español. Tampoco existe evidencia sobre qué componentes de
una arquitectura de ese tipo explican su comportamiento, ni sobre su efecto en el
trabajo del docente.
\end{estado}

Esta brecha es estrecha por diseño. No se afirma que falte investigación sobre IA en
educación, ni sobre RAG, ni sobre detección de contradicciones. Se afirma que la
intersección específica --- inconsistencia intracorpus como objeto, corpus de curso
como dominio, español como idioma, y atribución del desempeño a componentes --- no
está cubierta.

\subsection{Aporte que se pretende estudiar}
\label{sec:brecha-aporte}

El aporte que esta investigación persigue no es el sistema. Es el conocimiento que su
evaluación puede producir, en cuatro frentes:

\begin{enumerate}
  \item \textbf{Caracterización del desempeño por clase de defecto.} Determinar con
        qué precisión y \emph{recall} se detecta cada clase --- contradicción,
        redundancia, desactualización --- y con qué tasa de falsos positivos sobre
        controles negativos, reportando las clases por separado dada la asimetría de
        dificultad documentada en la línea C.
  \item \textbf{Atribución del desempeño a componentes.} Determinar cuánto aporta cada
        elemento de la arquitectura mediante configuraciones de comparación y
        ablaciones, en lugar de evaluar el sistema como un todo indivisible.
  \item \textbf{Relación entre confianza estimada y corrección observada.}
        Determinar si el puntaje que el sistema produce guarda correspondencia
        estadística con su tasa de acierto y, en caso negativo, si un método de
        calibración establecido la restablece.
  \item \textbf{Efecto sobre el trabajo del docente.} Determinar cómo se modifican el
        tiempo de ejecución y la carga percibida frente a una condición de revisión
        sin asistencia, midiendo ambas dimensiones por la advertencia de
        \textcite{su2026efficiency} sobre la insuficiencia del tiempo como único
        indicador.
\end{enumerate}

\subsection{Productos de conocimiento previstos}
\label{sec:brecha-productos}

Además de los resultados, la investigación prevé producir tres artefactos de
conocimiento reutilizables con independencia del sistema:

\begin{itemize}
  \item Un \textbf{benchmark anotado} de inconsistencias en documentación de cursos en
        español, con defectos plantados, controles negativos, evidencia esperada y
        acuerdo entre anotadores reportado.
  \item Un \textbf{protocolo de evaluación} que adapte los marcos de la línea D$'$ a
        un corpus en español, incluida la validación de los jueces automáticos contra
        anotación humana.
  \item Una \textbf{taxonomía operacional de defectos} de documentación de cursos,
        derivada de la literatura y verificable en la anotación.
\end{itemize}

\begin{advertencia}
Estos tres productos son \emph{previstos}, no existentes. El benchmark inicial
descrito en la Parte~IV es una versión piloto sin anotación independiente ni acuerdo
entre anotadores calculado. Presentarlo hoy como benchmark validado sería una
afirmación de la categoría (A) sin evidencia que la sustente.
\end{advertencia}
